% ============================================================
%  CAPÍTULO 5 — ARQUITECTURA DEL SISTEMA
% ============================================================
\chapter{Arquitectura del sistema}
\label{ch:arquitectura}

\section{Visión global}
\label{sec:vision-global}

El sistema se estructura en tres capas claramente diferenciadas
(figura~\ref{fig:arch-overview}):

\begin{enumerate}
  \item \textbf{Base de datos:} PostgreSQL gestionado por SQLAlchemy~2.x.
  \item \textbf{Backend (\textit{uniback}):} FastAPI + sistema de plugins; expone una
    \gls{api} REST que incluye endpoints \gls{crud} estándar, el endpoint de pantallas
    dinámicas y el endpoint de operaciones en bloque.
  \item \textbf{Frontend (\textit{ngt-gui}):} librería Angular que consume los esquemas
    y las pantallas dinámicas para renderizar la \gls{ui} sin lógica hardcoded.
\end{enumerate}

% TODO: insertar diagrama de arquitectura
\begin{figure}[H]
  \centering
  % \includegraphics[width=0.9\textwidth]{Media/Arquitectura/arch_overview.png}
  \fbox{\parbox{0.85\textwidth}{\centering\vspace{2cm}
    [Insertar diagrama de arquitectura general]\vspace{2cm}}}
  \caption{Visión general de la arquitectura del sistema.}
  \label{fig:arch-overview}
\end{figure}

\section{Tubería de transformación de esquemas}
\label{sec:pipeline}

La tubería de transformación es el núcleo funcional del \textit{framework}. Su objetivo es
convertir la definición de un modelo SQLAlchemy en un JSON~Schema anotado que el
\textit{frontend} puede consumir directamente.

\begin{figure}[H]
  \centering
  \fbox{\parbox{0.85\textwidth}{\centering\vspace{1.5cm}
    SQLAlchemy Model → Pydantic Schema → JSON Schema (draft-07) →
    x-* Annotations → Screen Definition\vspace{1.5cm}}}
  \caption{Etapas de la tubería de transformación de esquemas.}
  \label{fig:pipeline}
\end{figure}

Cada etapa añade información semántica:

\begin{description}
  \item[SQLAlchemy~→~Pydantic:] los metadatos de columna (\texttt{String}, \texttt{Integer},
    \texttt{ForeignKey}, \ldots) se convierten en tipos Pydantic con sus validadores.
  \item[Pydantic~→~JSON~Schema:] \texttt{model\_json\_schema()} produce un esquema estándar
    con \texttt{type}, \texttt{format}, \texttt{minLength}, \ldots
  \item[JSON~Schema~→~x-* annotations:] el sistema de plugins añade campos como
    \texttt{x-label}, \texttt{x-widget}, \texttt{x-foreign-key}, \texttt{x-options-endpoint}
    y \texttt{x-order}, que el \textit{frontend} interpreta.
  \item[Screen Definition:] el endpoint \texttt{/screens/<nombre>} ensambla el esquema
    anotado junto con metadatos de pantalla (\texttt{screen\_type}, \texttt{title},
    \texttt{actions}, \texttt{endpoint}) en un objeto que el \textit{frontend} consume
    directamente para renderizar la vista.
\end{description}

\section{Sistema de plugins por capas}
\label{sec:plugins}

\subsection{Motivación}
\label{subsec:plugins-motivacion}

La transformación de esquemas requiere enriquecimientos heterogéneos que varían según la
entidad: algunos campos necesitan un selector desplegable con opciones cargadas desde el
\gls{api}, otros requieren validadores especiales, y algunos deben ocultarse en ciertas
vistas. Centralizar toda esta lógica en un único generador lo haría inmanejable y difícil
de extender.

\subsection{Diseño}
\label{subsec:plugins-diseno}

Cada plugin es una clase Python con una interfaz común:

\begin{lstlisting}[language=Python, caption={Interfaz base de un plugin de esquema.}]
class SchemaPlugin(Protocol):
    priority: int  # menor = se aplica antes

    def apply(
        self,
        schema: dict,
        model: type[DeclarativeBase],
        field_name: str,
        column_info: MappedColumn,
    ) -> dict:
        """Devuelve el schema enriquecido."""
        ...
\end{lstlisting}

Los plugins se registran en un registro global ordenado por prioridad y se aplican en
cascada. El sistema soporta plugins de \textbf{columna} (operan campo a campo) y plugins
de \textbf{modelo} (operan sobre el esquema completo).

\subsection{Plugins implementados}
\label{subsec:plugins-implementados}

\begin{table}[H]
  \centering
  \caption{Plugins de esquema implementados en uniback.}
  \label{tab:plugins}
  \begin{tabular}{lp{9cm}}
    \toprule
    \textbf{Plugin} & \textbf{Función} \\
    \midrule
    \texttt{LabelPlugin} & Añade \texttt{x-label} a partir del nombre de columna
      (snake\_case → Title Case) o de metadatos \texttt{info}. \\
    \texttt{ForeignKeyPlugin} & Detecta claves foráneas y añade
      \texttt{x-foreign-key} y \texttt{x-options-endpoint} para carga de opciones. \\
    \texttt{ValidatorPlugin} & Traduce restricciones SQLAlchemy (\texttt{String(n)},
      \texttt{CheckConstraint}, \ldots) a \texttt{x-validators}. \\
    \texttt{WidgetPlugin} & Asigna el widget Formly más apropiado
      (\texttt{x-widget}: input, select, date, textarea, \ldots) según el tipo. \\
    \texttt{ReadonlyPlugin} & Marca campos autogenerados (\texttt{id}, \texttt{created\_at})
      como \texttt{x-readonly}. \\
    \texttt{OrderPlugin} & Asigna \texttt{x-order} según el orden de declaración en el
      modelo. \\
    \bottomrule
  \end{tabular}
\end{table}

\section{Arquitectura del frontend}
\label{sec:arch-frontend}

\subsection{Estructura de la librería ngt-gui}
\label{subsec:ngtgui-estructura}

\textit{ngt-gui} es una librería Angular construida con \textbf{ng-packagr} y distribuida
como paquete npm. Su estructura interna sigue la organización por funcionalidades:

\begin{verbatim}
projects/ngt-gui/src/lib/
  core/
    schema-to-formly.ts      # JSON Schema → Formly config
    schema-to-aggrid.ts      # JSON Schema → AG Grid ColDef
    entity-client.service.ts # Cliente HTTP genérico
  components/
    dynamics/
      dynamic-page/          # Enrutador de vistas dinámicas
      dynamic-browse/        # Tabla + filtros + botones CRUD
      dynamic-form/          # Formulario generado (Formly)
      bulk-edit-grid/        # Cuadrícula edición masiva
      formly-editor/         # Editor visual de esquemas
\end{verbatim}

\subsection{Flujo de datos en una vista dinámica}
\label{subsec:flujo-datos}

\begin{enumerate}
  \item El usuario navega a \texttt{/d/<screenName>}.
  \item \texttt{DynamicPageResolver} llama a \texttt{GET /screens/<screenName>} y
    almacena la definición en \texttt{ActivatedRoute.data}.
  \item \texttt{DynamicPageComponent} lee el \texttt{screen\_type} y renderiza el
    componente adecuado: \texttt{DynamicBrowseComponent}, \texttt{DynamicFormComponent}
    o \texttt{BulkEditGridComponent}.
  \item Cada componente usa el \texttt{entityPath} del objeto de pantalla para
    construir las URLs de datos (lista, detalle, bulk) usando el servicio
    \texttt{BACKEND\_SERVICE} como fuente de la URL base y las cabeceras de
    autenticación.
\end{enumerate}

\subsection{Contrato de pantalla dinámica}
\label{subsec:contrato-pantalla}

\begin{lstlisting}[language=json, caption={Ejemplo de objeto de pantalla dinámica (screen definition).}]
{
  "name": "identities_browser",
  "screen_type": "browser",
  "endpoint": "/api/identities",
  "definition": {
    "title": "Identidades",
    "fieldId": "id",
    "actions": ["create", "edit", "delete", "view"],
    "columns": [
      {
        "field": "name",
        "label": "Nombre",
        "type": "string",
        "editable": true,
        "required": true,
        "validators": { "maxLength": 255 }
      }
    ]
  },
  "entity_schema": { ... }
}
\end{lstlisting}

\section{Endpoint de operaciones en bloque}
\label{sec:bulk-endpoint}

El endpoint \texttt{POST /<entidad>/bulk} recibe un array de objetos con la siguiente
estructura:

\begin{lstlisting}[language=json, caption={Payload del endpoint /bulk.}]
[
  { "data": { "name": "Nueva fila" } },
  { "id": 42, "data": { "name": "Fila modificada" } },
  { "_op": "delete", "id": 17 }
]
\end{lstlisting}

La respuesta incluye el resultado de cada operación y un resumen:

\begin{lstlisting}[language=json, caption={Respuesta del endpoint /bulk.}]
{
  "results": [
    { "ok": true, "id": 101 },
    { "ok": true, "id": 42 },
    { "ok": false, "errors": [{ "loc": ["name"], "msg": "..." }] }
  ],
  "summary": { "created": 1, "updated": 1, "deleted": 0, "failed": 1 }
}
\end{lstlisting}

El parámetro de consulta \texttt{?on\_error=partial|abort} controla el comportamiento
transaccional:

\begin{itemize}
  \item \textbf{partial:} cada operación se ejecuta en su propio savepoint SQL; los fallos
    no afectan a las operaciones exitosas.
  \item \textbf{abort:} toda la petición se ejecuta en una única transacción; cualquier
    error provoca el rollback de todas las operaciones.
\end{itemize}

\section{Decisiones arquitectónicas relevantes}
\label{sec:decisiones}

\subsection{Standalone components en Angular}
\label{subsec:standalone}

Todos los componentes de \textit{ngt-gui} se declaran como \texttt{standalone: true},
eliminando la necesidad de módulos Angular. Esto simplifica la distribución de la librería
y reduce el riesgo de conflictos de módulos con las aplicaciones consumidoras.

\subsection{Inyección de dependencias con tokens opacos}
\label{subsec:di-tokens}

El servicio de \textit{backend} se inyecta mediante el token opaco \texttt{BACKEND\_SERVICE}
en lugar de una clase concreta. Esto permite que cualquier aplicación proporcione su propia
implementación del servicio sin modificar la librería, siguiendo el principio de inversión
de dependencias.

\subsection{HttpClient directo vs. EntityClient}
\label{subsec:httpclient}

El componente \texttt{BulkEditGridComponent} usa \texttt{HttpClient} directamente en lugar
del servicio \texttt{EntityClient} de la librería. Esta decisión evita un ciclo de
dependencias de inyección (\texttt{NG0200}) que se produce cuando un componente
\textit{standalone} de la librería intenta inyectar un servicio que, a su vez, depende de
la inicialización de la aplicación.
